\section{Results and Ablations}
\label{sec:results}

The assembled detector reaches $F_1=\num{0.931}$ on \PCB{}
($P=\num{1.000}$, $R=\num{0.871}$, 14 components) and $F_1=\num{1.000}$ on the
synthetic \CRACK{} sample. NMS collapses the gradient shoulder to a 1-px ridge
and hysteresis isolates the crack from background texture; the full per-stage
gallery (grayscale $\to$ smoothing $\to$ $|M|$ $\to$ orientation $\to$ NMS $\to$
strong/weak $\to$ hysteresis) for both benchmarks is in the supplementary.

\ifdefined\SUPPLEMENTARY
\begin{figure}[tb]
  \centering
  \includegraphics[width=\linewidth]{task3_canny_concrete.png}
  \caption{Four-stage Canny on a concrete-crack sample: grayscale, Gaussian
  pre-smoothing, Sobel $|M|$, orientation (hue)$\times|M|$ (value), NMS,
  double-threshold classes (weak = blue, strong = yellow), hysteresis, ground
  truth.}
  \label{fig:pipeline}
\end{figure}
\fi

\subsection{Ablation 1 -- noise and Gaussian pre-smoothing}
Differentiating a noisy image amplifies the noise: the gradient-magnitude std.\
dev.\ grows with $\sigma_n$ as $3.50\!\to\!7.30\!\to\!11.42$ for
$\sigma_n\in\{10,25,40\}$. A single Gaussian pass ($\sigma=1.4$) holds it at
$1.47\!\to\!1.79\!\to\!2.42$ --- a noise-suppression ratio rising from
\SI{2.4}{\times} to \SI{4.7}{\times} --- while keeping the crack ridge intact
(\cref{fig:abl1}). This is the derivative-of-noise argument for Stage 1.

\begin{figure}[tb]
  \centering
  \ifdefined\SUPPLEMENTARY
    \includegraphics[width=0.86\linewidth]{task4_ablation1_noise_smoothing_full.png}
    \caption{Ablation 1. Rows: $\sigma_n\in\{10,25,40\}$. Columns: noisy input;
    $|\nabla|$ without pre-smoothing; $|\nabla|$ with Gaussian pre-smoothing.}
  \else
    \includegraphics[width=0.82\linewidth]{task4_ablation1_noise_smoothing.png}
    \caption{Ablation 1 (extremes $\sigma_n\in\{10,40\}$; $\sigma_n{=}25$ in the
    supplementary). Columns: noisy input; $|\nabla|$ without pre-smoothing;
    $|\nabla|$ with Gaussian pre-smoothing.}
  \fi
  \label{fig:abl1}
\end{figure}

\subsection{Ablation 2 -- non-maximum suppression}
Thresholding the (over-smoothed) gradient magnitude directly yields thick bands
covering \SI{40.0}{\percent} of the frame; applying the same threshold
\emph{after} NMS yields 1-px ridges at \SI{4.9}{\percent} density --- a
\SI{8.1}{\times} thinning (\cref{fig:abl2}).

\begin{figure}[tb]
  \centering
  \includegraphics[width=0.9\linewidth]{task4_ablation2_nms.png}
  \caption{Ablation 2. $|M|\!\ge\!t$ before NMS (thick bands) vs.\
  $\mathrm{NMS}(|M|)\!\ge\!t$ (1-px ridges), same threshold $t$.}
  \label{fig:abl2}
\end{figure}

\subsection{Ablation 3 -- hysteresis vs.\ single threshold}
On a faint, contrast-varying crack in noise (\cref{fig:abl3},
\cref{tab:ablation3}): a single low threshold floods the background
($P=\num{0.535}$, 225 fragments); a single high threshold is clean but breaks
the crack ($R=\num{0.898}$, 40 fragments); double-threshold hysteresis links
the weak crack stretches to strong seeds while rejecting isolated texture,
reaching $F_1=\num{0.988}$ in only \textbf{12} connected components.

\begin{table}[t]
  \centering
  \caption{Ablation 3: single threshold vs.\ hysteresis (faint crack, $\sigma_n{=}12$).}
  \label{tab:ablation3}
  \small
  \setlength{\tabcolsep}{4pt}
  \begin{tabular}{lccccc}
    \toprule
    method & $P$ & $R$ & $F_1$ & density & \#comp. \\
    \midrule
    single @ $\Tlo$  & 0.535 & 1.000 & 0.697 & 0.019 & 225 \\
    single @ $\Thi$  & 0.993 & 0.898 & 0.943 & 0.006 & 40 \\
    \textbf{hysteresis} & 0.982 & 0.994 & \textbf{0.988} & 0.010 & \textbf{12} \\
    \bottomrule
  \end{tabular}
\end{table}

\begin{figure}[tb]
  \centering
  \includegraphics[width=0.92\linewidth]{task4_ablation3_hysteresis.png}
  \caption{Ablation 3. Left to right: noisy input; single threshold at $\Tlo$;
  single threshold at $\Thi$; double-threshold hysteresis.}
  \label{fig:abl3}
\end{figure}

\subsection{Scale-space causality}
\label{sec:scalespace}
Running the detector at $\sigma\in\{1.0,2.5,5.0\}$ gives edge counts
$2055 \to 1418 \to 881$ --- strictly decreasing (\cref{fig:scale}). Tracking
the second-derivative zero-crossings along a fixed scan-line over
$\sigma\in[0.4,8.0]$ likewise shows the count fall monotonically from
\num{120} to \num{15}.

\emph{Why edges only merge, never appear (Witkin
causality~\cite{witkin1983}).} Gaussian smoothing of a 1-D
signal is the heat equation $\partial_t u = \tfrac12\partial_{xx}u$ run for
time $t=\sigma^2$. Diffusion obeys a maximum principle: it cannot increase the
number of local extrema of $u$ or of any derivative $\partial_x^n u$. Edges are
zero-crossings of $\partial_{xx}(G_\sigma\conv f)$, so no new edge can be
created; existing zero-crossings drift and annihilate in pairs as neighbouring
extrema of $\partial_{xx}u$ merge. Every coarse-scale edge therefore has a
unique fine-scale ancestor, which is what makes coarse-to-fine tracking well
posed. (Full argument: supplementary, Appendix~A.)

\begin{figure}[tb]
  \centering
  \ifdefined\SUPPLEMENTARY
    \includegraphics[width=\linewidth]{task5_scale_space_full.png}
  \else
    \includegraphics[width=\linewidth]{task5_scale_space.png}
  \fi
  \caption{Scale-space: increasing $\sigma$ destroys texture edges; the crack
  edges shift and merge but no new edges appear. Edge counts
  $2055 \to 1418 \to 881$.}
  \label{fig:scale}
\end{figure}

\ifdefined\SUPPLEMENTARY
\begin{figure}[tb]
  \centering
  \includegraphics[width=\linewidth]{task5_scale_space_causality.png}
  \caption{Left: zero-crossing fingerprint (branches only merge upward).
  Right: zero-crossing count is monotone non-increasing in $\sigma$.}
  \label{fig:causality}
\end{figure}
\fi

\subsection{Parameter sensitivity}
\label{sec:sensitivity}
We sweep $\Thi$ (quantiles $0.75/0.88/0.96$ of the NMS response) against
$\Tlo=\{0.3,0.5,0.8\}\,\Thi$ on noisy \CRACK{} and \PCB{} (full $3\times3$
edge-map grids in \cref{fig:grid}). \Cref{tab:grid} extracts the diagonal: a low
$\Thi$ admits texture regardless of $\Tlo$ ($F_1\!\approx\!0.15$); the usable
regime is a high $\Thi$ (top few percent of gradient energy) with $\Tlo$ at
\SIrange{50}{80}{\percent} of it ($F_1=\num{0.888}$), matching the classic
$2$--$3\!:\!1$ ratio heuristic. The \PCB{} grid peaks at $F_1=\num{0.936}$
($\Thi=\text{q}0.88$, $\Tlo=0.3\,\Thi$), reflecting its higher edge contrast.

\begin{table}[t]
  \centering
  \caption{Sensitivity diagonal, noisy \CRACK{} ($\sigma_n{=}25$); best in bold.
  Full $3\times3$ grids: \cref{fig:grid}.}
  \label{tab:grid}
  \small
  \setlength{\tabcolsep}{4.5pt}
  \begin{tabular}{llccc}
    \toprule
    $\Thi$ (quantile) & $\Tlo/\Thi$ & density & \#comp. & $F_1$ \\
    \midrule
    q0.75 & 0.30 & 0.291 & 227 & 0.151 \\
    q0.88 & 0.50 & 0.148 & 270 & 0.268 \\
    q0.96 & 0.30 & 0.082 & 39  & 0.436 \\
    q0.96 & 0.50 & 0.035 & 47  & 0.769 \\
    q0.96 & 0.80 & 0.020 & 77  & \textbf{0.888} \\
    \bottomrule
  \end{tabular}
\end{table}

\ifdefined\SUPPLEMENTARY
\begin{figure}[tb]
  \centering
  \includegraphics[width=\linewidth]{task5_sensitivity_grid_concrete.png}
  \includegraphics[width=\linewidth]{task5_sensitivity_grid_pcb.png}
  \caption{$3\times3$ $(\Thi,\Tlo)$ grids on noisy concrete (top) and PCB
  (bottom). Rows: $\Thi$ quantile increasing downward. Columns: $\Tlo/\Thi$
  increasing rightward.}
  \label{fig:grid}
\end{figure}
\else
\begin{figure}[tb]
  \centering
  \includegraphics[width=0.9\linewidth]{task5_sensitivity_grid_concrete.png}
  \caption{$3\times3$ $(\Thi,\Tlo)$ grid on noisy concrete. Rows: $\Thi$
  quantile $\uparrow$ downward; columns: $\Tlo/\Thi$ $\uparrow$ rightward.
  PCB grid in the supplementary.}
  \label{fig:grid}
\end{figure}
\fi
